\section{Evaluation}
\label{sec:evaluation}

\citet{mobilitydt2023} separate \emph{fidelity} (precision independent of latency) from
\emph{physical-to-digital synchronicity} (discrepancy caused by latency); the separation
is the useful move, because the two trade against each other and one ``quality'' number
hides the trade. Their own measurement makes the point: compressing offloaded imagery
from $>18$~MB to $\approx6$~MB cost less than 0.1\% of PSNR while cutting latency
substantially. The analogues here are \emph{encoding fidelity} (rendered property versus
the encoding evaluated in the pure layer --- an automatable readback test) and
\emph{synchronicity} (p50/p95/p99 publish-to-pixel latency at a stated element count).

Three further measures follow from the new substrates and from \tech{9}.
\textbf{Residual credibility}: report the residual's own error budget beside it, since a
ghost twin showing 2\,\textdegree C divergence is useless if pipeline uncertainty is
3\,\textdegree C --- \citet{liao2025digital} model this discipline with quantitative
metrics rather than visual plausibility. \textbf{Registration accuracy}: homography
reprojection error, expressed in ground-plane metres as well as pixels, because three
pixels near the horizon of a traffic photograph can be several metres.
\textbf{Synchronisation skew}: the signed video-frame-to-playhead difference, reported
continuously since transport latency drifts; the question is not whether skew is zero
but whether it is known, bounded and smaller than the phenomenon being compared. Finally,
\citet{chen2024tangible} evaluate an MR airport digital tower with ten licensed
controllers on workload, situational awareness and trust --- the right constructs for an
operator-facing view, and not substitutable by usability, since a view can be easy to use
and still degrade situational awareness by drawing attention to the wrong thing. All but
the operator study belong in CI from the phase that introduces them.
